\chapter{Reading GRAIL Output}
\label{ch:output}
GRAIL returns a displacement for the residual and derivative callbacks supplied
by the caller. Solver convergence and accuracy of the physical prediction are
separate checks.

\section{The Result Dataclass}
\label{sec:result-overview}
\code{ResolventResult} contains \code{delta\_x\_star},
\code{iter\_count}, \code{residual\_history},
\code{condition\_estimate}, \code{pca\_rank}, and
\code{conformal\_intervals}. Arrays use the selected backend.
Use \code{np.asarray(result.delta\_x\_star)} for a portable NumPy copy.

\section{The Predicted Shift}
\label{sec:delta-x-star}
With the emulator-minus-reference residual
$b=\mathcal T_e(x_e)-\mathcal T_t(x_e)$, the returned displacement
estimates $x_e-x_t$. Therefore
\begin{equation}
 x_t^{\mathrm{estimate}}=x_e-\delta x^{*,\mathrm{pred}}.
\end{equation}
\begin{lstlisting}[style=grailInStyle]
reference_estimate = x_emulator - result.delta_x_star
\end{lstlisting}
Keep that residual sign convention in both the callback wiring and subsequent
correction. The approximation concerns a small smooth perturbation, while
repeated linear feedback is retained to all orders.

\section{GMRES Diagnostics}
\label{sec:gmres-diag}
\subsection{Iteration Count}
\label{sec:iter-count}
\code{iter\_count} records the number of GMRES iterations. Stopping before
the maximum is not a convergence certificate: early breakdown can also stop an
iteration. Check the residual history and independently evaluate the equation
with the returned displacement.
\subsection{Residual History and True Equation Defect}
\label{sec:residual-history}
For $p=\code{result.delta\_x\_star}$, recompute
$r=p-\mathrm{emulator\_jvp}(\mathrm{host\_jvp}(p))-b$.
Compare $\|r\|/\|b\|$ with the requested tolerance when $b\ne0$; for a
zero right-hand side use an absolute check. This tests the supplied operator.
Derivative verification against the actual host cycle is an additional step.
\subsection{Condition Estimate}
\label{sec:condition}
\code{condition\_estimate} is obtained from a projected Arnoldi/Hessenberg
problem. It is a diagnostic, not a certificate of the full-space condition
number or physical reliability. An infinite estimate can reflect rank deficiency,
breakdown, or a singular projected problem. Interpret it together with the
true residual, independent derivative checks, and physical validation.

\section{PCA Rank}
\label{sec:pca-rank}
A bare \code{resolvent\_solve} result has \code{pca\_rank=0}.
PCA residual representation is a separate layer. A zero rank field does not
establish whether a caller calibrated a prediction elsewhere.

\section{Conformal Bands}
\label{sec:conformal}
A bare solve returns \code{conformal\_intervals=None}; it does not construct
a band automatically. Calibration needs held-out prediction errors against the
chosen reference target, with a score and sampling protocol fixed beforehand.
Chapter~\ref{ch:coverage} gives the exact finite-sample procedure.

\section{A Complete Worked Example}
\label{sec:worked-complete}
\begin{lstlisting}[style=grailInStyle]
python examples/structured_response.py
\end{lstlisting}
This example uses the current learned neural derivative at the surrogate
equilibrium, reports its true linear defect, and compares the corrected
population profile against an independently solved reference. The reported
shift-prediction error is relative to the actual shift; it is distinct from
population-relative error or the magnitude of amplification itself.
